% Section 7 — First-proposal performance across games (SYNTHETIC results)
\section{First-proposal performance}
\label{sec:results}

\Cref{tab:main,fig:primary} give the primary result in each game \synbadge. Every number in this section
is synthetic and illustrates the evidence the confirmation runs must produce.

\begin{table}[t]
\caption{Primary results on sealed families: failure-inclusive utility of one committed proposal
against the locked strongest non-search control, with 95\% paired family-bootstrap intervals, and the
search budget whose quality one proposal matches. All values synthetic (gaps EZ-G01 to EZ-G06).}
\label{tab:main}
\centering\tablesf\setlength\tabcolsep{4.6pt}
\begin{tabular}{@{}llrrrcrr@{}}
\toprule
 & & \multicolumn{2}{c}{Utility} & & & & \\
\cmidrule(lr){3-4}
Game (track) & Locked comparator & EZ & Comp. & $\Delta$ & 95\% interval & Failed & $B_{\mathrm{eq}}$ \\
\midrule
A Learning procedures (C0) & tree selector & \syn[EZ-G01]{\EZuA} & \syn[EZ-G01]{\EZcA} & \syn[EZ-G01]{\EZdA} & \syn[EZ-G01]{[\EZdAlo, \EZdAhi]} & \syn[EZ-G01]{\EZfailA\%} & \syn[EZ-G02]{\EZbeqA} \\
B1 Next experiment (CH) & Gaussian-process EI & \syn[EZ-G03]{\EZuB} & \syn[EZ-G03]{\EZcB} & \syn[EZ-G03]{\EZdB} & \syn[EZ-G03]{[\EZdBlo, \EZdBhi]} & \syn[EZ-G03]{\EZfailB\%} & --- \\
C CAD parts (C0) & retrieval + adaptation & \syn[EZ-G05]{\EZuC} & \syn[EZ-G05]{\EZcC} & \syn[EZ-G05]{\EZdC} & \syn[EZ-G05]{[\EZdClo, \EZdChi]} & \syn[EZ-G05]{\EZfailC\%} & \syn[EZ-G06]{\EZbeqC} \\
\bottomrule
\end{tabular}
\eztablenote{The tree selector chooses among complete archived procedures;
Gaussian-process EI sees the same frozen history; retrieval adapts the nearest training part by a
fixed rule. $B_{\mathrm{eq}}$ counts evaluated candidates (A) or simulations (C); B1 is a single
decision and has no search-equivalent budget.}
\end{table}

\begin{figure}[t]
\centering
% Figure — first-proposal primary results: family-level paired differences per game.
% Reads SYNTHETIC data (gaps EZ-G01, EZ-G03, EZ-G05).
\pgfplotstableread[col sep=comma]{data/synthetic/primary.csv}\ezprimary
\pgfplotstablegetelem{0}{delta}\of\ezprimary\edef\ezdA{\pgfplotsretval}
\pgfplotstablegetelem{0}{lo}\of\ezprimary\edef\ezloA{\pgfplotsretval}
\pgfplotstablegetelem{0}{hi}\of\ezprimary\edef\ezhiA{\pgfplotsretval}
\pgfplotstablegetelem{1}{delta}\of\ezprimary\edef\ezdB{\pgfplotsretval}
\pgfplotstablegetelem{1}{lo}\of\ezprimary\edef\ezloB{\pgfplotsretval}
\pgfplotstablegetelem{1}{hi}\of\ezprimary\edef\ezhiB{\pgfplotsretval}
\pgfplotstablegetelem{2}{delta}\of\ezprimary\edef\ezdC{\pgfplotsretval}
\pgfplotstablegetelem{2}{lo}\of\ezprimary\edef\ezloC{\pgfplotsretval}
\pgfplotstablegetelem{2}{hi}\of\ezprimary\edef\ezhiC{\pgfplotsretval}
\begin{tikzpicture}
\begin{groupplot}[ez,
    group style={group size=3 by 1, horizontal sep=0.45cm, y descriptions at=edge left},
    width=0.385\linewidth, height=4.6cm, ymin=0, ymax=62,
    xlabel={paired utility difference}, ylabel={families},
    xtick={-0.4,0,0.4}, ytick={0,20,40,60}, xticklabel style={/pgf/number format/fixed},
    xmin=-0.62, xmax=0.68, axis lines*=left, clip=false,
    ezhist/.style={hist={bins=26, data min=-0.62, data max=0.68}, fill=ezaccent!32, draw=ezaccent!75, line width=0.3pt}]
  \nextgroupplot[title={A \textperiodcentered{} learning procedures}]
    \fill[ezwarm, opacity=0.16] (axis cs:\ezloA,0) rectangle (axis cs:\ezhiA,62);
    \addplot[ezhist] table[col sep=comma, y=delta] {data/synthetic/family-A.csv};
    \draw[ezgrey, line width=0.5pt, densely dashed] (axis cs:0,0) -- (axis cs:0,62);
    \draw[ezwarm, line width=1.1pt] (axis cs:\ezdA,0) -- (axis cs:\ezdA,62);
  \nextgroupplot[title={B1 \textperiodcentered{} next experiment}]
    \fill[ezwarm, opacity=0.16] (axis cs:\ezloB,0) rectangle (axis cs:\ezhiB,62);
    \addplot[ezhist] table[col sep=comma, y=delta] {data/synthetic/family-B.csv};
    \draw[ezgrey, line width=0.5pt, densely dashed] (axis cs:0,0) -- (axis cs:0,62);
    \draw[ezwarm, line width=1.1pt] (axis cs:\ezdB,0) -- (axis cs:\ezdB,62);
  \nextgroupplot[title={C \textperiodcentered{} CAD parts}]
    \fill[ezwarm, opacity=0.16] (axis cs:\ezloC,0) rectangle (axis cs:\ezhiC,62);
    \addplot[ezhist] table[col sep=comma, y=delta] {data/synthetic/family-C.csv};
    \draw[ezgrey, line width=0.5pt, densely dashed] (axis cs:0,0) -- (axis cs:0,62);
    \draw[ezwarm, line width=1.1pt] (axis cs:\ezdC,0) -- (axis cs:\ezdC,62);
\end{groupplot}
\synfigtag
\end{tikzpicture}

\caption{Family-level paired differences between one EurekaZero proposal and the locked comparator.
The orange line is the family-weighted mean and the shaded band its 95\% interval; the dashed line
marks no difference. Failed first proposals score $-1$, are collected in the leftmost bin and stay in
every mean. Synthetic (gaps EZ-G01, EZ-G03, EZ-G05).}
\label{fig:primary}
\end{figure}

\subsection{Learning procedures}

On \syn[EZ-G20]{\EZnA} sealed families, one committed procedure improves failure-inclusive utility over
the tree selector by \syn[EZ-G01]{\EZdA} [\syn[EZ-G01]{\EZdAlo}, \syn[EZ-G01]{\EZdAhi}], with
\syn[EZ-G01]{\EZfailA\%} of first proposals failing. The gain is not uniform (\cref{fig:primary}):
it is concentrated in small and wide tables, where the archive's best procedure for the nearest
family transfers poorly, and it vanishes on large, low-noise regression tasks, where every strong
method converges (\cref{sec:heterogeneity}). Against TabPFN, whose output is a prediction rather than
a procedure, EurekaZero's held-out error ratio is \syn[EZ-G08]{\EZtabpfnRatio} [\syn[EZ-G08]{\EZtabpfnLo},
\syn[EZ-G08]{\EZtabpfnHi}]: comparable accuracy from a standalone program that fits in a median of
\syn[EZ-G08]{\EZfitMedian} seconds and runs without the model that wrote it.

\subsection{Next-experiment decisions}

Given identical frozen histories (B1), one EurekaZero query improves on the next point chosen by
Gaussian-process expected improvement \citep{snoek2012bo,balandat2020botorch} by \syn[EZ-G03]{\EZdB}
[\syn[EZ-G03]{\EZdBlo}, \syn[EZ-G03]{\EZdBhi}], and \syn[EZ-G03]{\EZvalidB\%} of its queries are valid.
Closed loop (B2, \cref{fig:curves}), EurekaZero reaches the target regret of 0.15 in
\syn[EZ-G04]{\EZreachEZ} evaluations, against \syn[EZ-G04]{\EZreachZS} for a ZeroShotOpt-class pretrained
optimizer \citep{meindl2025zeroshotopt}, \syn[EZ-G04]{\EZreachOF} for an OptFormer-class model
\citep{chen2022optformer} and \syn[EZ-G04]{\EZreachGPEI} for Gaussian-process EI; after 20 evaluations
its regret equals what Gaussian-process EI reaches after \syn[EZ-G04]{\EZkeqBtwenty}, a factor of
\syn[EZ-G04]{\EZsaveB} in experiments. The closed-loop result is a statement about sequential
optimization at a fixed budget, not about solving an objective in one call.

\begin{figure}[t]
\centering
% Figure — what one proposal is worth: search budget curves (A, C) and closed-loop regret (B2).
% Reads SYNTHETIC data (gaps EZ-G02, EZ-G04, EZ-G06, EZ-G07).
\pgfplotstableread[col sep=comma]{data/synthetic/curve-lines.csv}\ezlines
\pgfplotstablegetelem{0}{ez}\of\ezlines\edef\ezA{\pgfplotsretval}
\pgfplotstablegetelem{0}{ez_lo}\of\ezlines\edef\ezAlo{\pgfplotsretval}
\pgfplotstablegetelem{0}{ez_hi}\of\ezlines\edef\ezAhi{\pgfplotsretval}
\pgfplotstablegetelem{0}{comparator}\of\ezlines\edef\cmpA{\pgfplotsretval}
\pgfplotstablegetelem{0}{retrieval}\of\ezlines\edef\retA{\pgfplotsretval}
\pgfplotstablegetelem{0}{frontier}\of\ezlines\edef\frA{\pgfplotsretval}
\pgfplotstablegetelem{1}{ez}\of\ezlines\edef\ezC{\pgfplotsretval}
\pgfplotstablegetelem{1}{ez_lo}\of\ezlines\edef\ezClo{\pgfplotsretval}
\pgfplotstablegetelem{1}{ez_hi}\of\ezlines\edef\ezChi{\pgfplotsretval}
\pgfplotstablegetelem{1}{comparator}\of\ezlines\edef\cmpC{\pgfplotsretval}
\pgfplotstablegetelem{1}{frontier}\of\ezlines\edef\frC{\pgfplotsretval}
\begin{tikzpicture}
\begin{groupplot}[ez,
    group style={group size=3 by 1, horizontal sep=1.35cm},
    width=0.34\linewidth, height=4.9cm, axis lines*=left, clip=false,
    legend style={font=\sffamily\fontsize{5.6pt}{6.4pt}\selectfont, draw=none, fill=none,
      cells={anchor=west}, row sep=-1.5pt, at={(0.5,-0.3)}, anchor=north}]
  % ---------------- A: utility against search budget
  \nextgroupplot[title={A \textperiodcentered{} learning procedures}, xmode=log, log basis x=2,
      xmin=0.8, xmax=320, ymin=0.26, ymax=0.47, xtick={1,4,16,64,256}, xticklabels={1,4,16,64,256},
      xlabel={candidates evaluated by search}, ylabel={mean utility},
      yticklabel style={/pgf/number format/fixed, /pgf/number format/precision=2},
      ]
    \addplot[name path=hiA, draw=none, forget plot] table[col sep=comma, x=k, y=hi] {data/synthetic/curve-A.csv};
    \addplot[name path=loA, draw=none, forget plot] table[col sep=comma, x=k, y=lo] {data/synthetic/curve-A.csv};
    \addplot[ezgrey!25, forget plot] fill between[of=hiA and loA];
    \addplot[ezgrey, line width=0.9pt, mark=*, mark size=1.3pt] table[col sep=comma, x=k, y=mean] {data/synthetic/curve-A.csv};
    \addlegendentry{search at budget $k$}
    \fill[ezwarm, opacity=0.14] (axis cs:0.8,\ezAlo) rectangle (axis cs:320,\ezAhi);
    \addplot[ezwarm, line width=1.2pt, domain=0.8:320, samples=2] {\ezA};
    \addlegendentry{EurekaZero, one proposal}
    \addplot[ezaccent, line width=0.7pt, densely dashed, domain=0.8:320, samples=2] {\cmpA};
    \addlegendentry{tree selector (comparator)}
    \addplot[ezaccent!60, line width=0.7pt, dotted, domain=0.8:320, samples=2] {\retA};
    \addlegendentry{retrieval}
    \addplot[ezgreen, line width=0.7pt, dashdotted, domain=0.8:320, samples=2] {\frA};
    \addlegendentry{best frontier model, $K=1$}
    \draw[ezwarm, line width=0.5pt, -{Stealth[length=1.5mm]}] (axis cs:\EZbeqA,0.448) -- (axis cs:\EZbeqA,\ezA);
    \node[font=\sffamily\fontsize{5.8pt}{6.4pt}\selectfont, text=ezwarm, anchor=south] at (axis cs:\EZbeqA,0.448) {$B_{\mathrm{eq}}=\EZbeqA$};
  % ---------------- B2: closed-loop regret
  \nextgroupplot[title={B2 \textperiodcentered{} closed-loop optimization}, ymode=log, xmin=0, xmax=50,
      ymin=0.03, ymax=1.0, xlabel={evaluations}, ylabel={median normalized regret},
      ytick={0.05,0.1,0.2,0.5,1}, yticklabels={0.05,0.1,0.2,0.5,1}, legend columns=2, legend style={column sep=3pt}]
    \addplot[ezwarm, line width=1.2pt] table[col sep=comma, x=t, y=ez] {data/synthetic/regret-B2.csv};
    \addlegendentry{EurekaZero (SEQ)}
    \addplot[ezgreen, line width=0.7pt] table[col sep=comma, x=t, y=zeroshot] {data/synthetic/regret-B2.csv};
    \addlegendentry{ZeroShotOpt-class}
    \addplot[ezgreen!55, line width=0.7pt, densely dashed] table[col sep=comma, x=t, y=optformer] {data/synthetic/regret-B2.csv};
    \addlegendentry{OptFormer-class}
    \addplot[ezaccent, line width=0.8pt] table[col sep=comma, x=t, y=gpei] {data/synthetic/regret-B2.csv};
    \addlegendentry{Gaussian-process EI}
    \addplot[ezaccent!55, line width=0.7pt, densely dashed] table[col sep=comma, x=t, y=tpe] {data/synthetic/regret-B2.csv};
    \addlegendentry{TPE}
    \addplot[ezamber, line width=0.7pt, densely dotted] table[col sep=comma, x=t, y=cmaes] {data/synthetic/regret-B2.csv};
    \addlegendentry{CMA-ES}
    \addplot[ezgrey, line width=0.7pt] table[col sep=comma, x=t, y=random] {data/synthetic/regret-B2.csv};
    \addlegendentry{random}
    \draw[ezrule, line width=0.4pt, densely dashed] (axis cs:0,0.15) -- (axis cs:50,0.15);
  % ---------------- C: utility against simulation budget
  \nextgroupplot[title={C \textperiodcentered{} CAD parts}, xmode=log, log basis x=2,
      xmin=0.8, xmax=320, ymin=-0.12, ymax=0.32, xtick={1,4,16,64,256}, xticklabels={1,4,16,64,256},
      xlabel={simulations spent by search}, ylabel={mean utility},
      yticklabel style={/pgf/number format/fixed, /pgf/number format/precision=1}]
    \addplot[name path=hiC, draw=none, forget plot] table[col sep=comma, x=k, y=hi] {data/synthetic/curve-C.csv};
    \addplot[name path=loC, draw=none, forget plot] table[col sep=comma, x=k, y=lo] {data/synthetic/curve-C.csv};
    \addplot[ezgrey!25, forget plot] fill between[of=hiC and loC];
    \addplot[ezgrey, line width=0.9pt, mark=*, mark size=1.3pt] table[col sep=comma, x=k, y=mean] {data/synthetic/curve-C.csv};
    \addlegendentry{topology optimization at $k$}
    \fill[ezwarm, opacity=0.14] (axis cs:0.8,\ezClo) rectangle (axis cs:320,\ezChi);
    \addplot[ezwarm, line width=1.2pt, domain=0.8:320, samples=2] {\ezC};
    \addlegendentry{EurekaZero, one proposal}
    \addplot[ezaccent, line width=0.7pt, densely dashed, domain=0.8:320, samples=2] {\cmpC};
    \addlegendentry{retrieval + adaptation}
    \addplot[ezgreen, line width=0.7pt, dashdotted, domain=0.8:320, samples=2] {\frC};
    \addlegendentry{best frontier model, $K=1$}
    \draw[ezwarm, line width=0.5pt, -{Stealth[length=1.5mm]}] (axis cs:\EZbeqC,0.29) -- (axis cs:\EZbeqC,\ezC);
    \node[font=\sffamily\fontsize{5.8pt}{6.4pt}\selectfont, text=ezwarm, anchor=south] at (axis cs:\EZbeqC,0.29) {$B_{\mathrm{eq}}=\EZbeqC$};
\end{groupplot}
\synfigtag
\end{tikzpicture}

\caption{What one proposal is worth. Left and right: mean utility of search at budget $k$ (grey,
95\% band) against one EurekaZero proposal (orange band); the arrow marks the search-equivalent budget
$B_{\mathrm{eq}}$, the smallest measured budget at which search matches the first proposal. Centre:
median normalized regret over 50 closed-loop evaluations from identical initial observations; the
dashed line is the target regret. Synthetic (gaps EZ-G02, EZ-G04, EZ-G06, EZ-G07).}
\label{fig:curves}
\end{figure}

\subsection{Engineering parts}

\syn[EZ-G05]{\EZbuildC\%} of committed construction programs build valid solids,
\syn[EZ-G05]{\EZcomplyC\%} satisfy every interface, envelope and process rule, and
\syn[EZ-G05]{\EZfeasC\%} also meet the stiffness and stress limits under converged analysis. Over
\syn[EZ-G20]{\EZnC} sealed specification families, utility improves on retrieval with a fixed adaptation
rule by \syn[EZ-G05]{\EZdC} [\syn[EZ-G05]{\EZdClo}, \syn[EZ-G05]{\EZdChi}]; feasible parts weigh
\syn[EZ-G19]{\EZmassC} [\syn[EZ-G19]{\EZmassClo}, \syn[EZ-G19]{\EZmassChi}] of the tuned template at
equal compliance (\cref{fig:cad}), and one proposal matches topology optimization after
\syn[EZ-G06]{\EZbeqC} simulations with a declared CAD conversion. Of
\syn[EZ-G17]{\EZphysN} prospectively selected parts that were fabricated and load-tested by a laboratory
not involved in their design, all met the load requirement, and measured stiffness differed from the
prediction by \syn[EZ-G17]{\EZphysErr\%} on average.

\begin{figure}[t]
\centering
% Figure — CAD: relative mass against relative compliance, and an illustrative committed part.
% Reads SYNTHETIC data (gap EZ-G19); the part drawing is a schematic, not an output.
\begin{tikzpicture}
\begin{axis}[ez, width=0.56\linewidth, height=5.4cm, xmin=0.62, xmax=1.42, ymin=0.6, ymax=1.5,
    xlabel={compliance relative to the specification limit}, ylabel={mass relative to tuned template},
    axis lines*=left, legend columns=2, legend style={at={(0.5,0.99)}, anchor=north, column sep=4pt, fill=none}, unbounded coords=discard,
    xticklabel style={/pgf/number format/fixed, /pgf/number format/precision=1},
    yticklabel style={/pgf/number format/fixed, /pgf/number format/precision=1},
    legend style={font=\sffamily\fontsize{5.6pt}{6.4pt}\selectfont, draw=none, fill=white, fill opacity=0.85, text opacity=1, cells={anchor=west}}]
  \fill[ezwarm, opacity=0.06] (axis cs:1,0.6) rectangle (axis cs:1.42,1.5);
  \node[font=\sffamily\fontsize{5.8pt}{6.4pt}\selectfont, text=ezwarm, anchor=north east] at (axis cs:1.41,0.68) {infeasible};
  \draw[ezwarm, line width=0.6pt, densely dashed] (axis cs:1,0.6) -- (axis cs:1,1.5);
  \addplot[only marks, mark=*, mark size=1.4pt, ezwarm, fill opacity=0.8, ezwhere={method}{ez}] table[col sep=comma, x=compliance, y=mass] {data/synthetic/cad-pareto.csv};
  \addlegendentry{EurekaZero, one proposal}
  \addplot[only marks, mark=square*, mark size=1.2pt, ezaccent, fill opacity=0.7, ezwhere={method}{topopt}] table[col sep=comma, x=compliance, y=mass] {data/synthetic/cad-pareto.csv};
  \addlegendentry{topology optimization, 32 simulations}
  \addplot[only marks, mark=triangle*, mark size=1.5pt, ezgrey, fill opacity=0.7, ezwhere={method}{template}] table[col sep=comma, x=compliance, y=mass] {data/synthetic/cad-pareto.csv};
  \addlegendentry{tuned template}
  \addplot[only marks, mark=diamond*, mark size=1.5pt, ezaccent!50, fill opacity=0.7, ezwhere={method}{retrieval}] table[col sep=comma, x=compliance, y=mass] {data/synthetic/cad-pareto.csv};
  \addlegendentry{retrieval + adaptation}
  \addplot[only marks, mark=o, mark size=1.4pt, ezgreen, ezwhere={method}{frontier}] table[col sep=comma, x=compliance, y=mass] {data/synthetic/cad-pareto.csv};
  \addlegendentry{best frontier model}
\end{axis}
% schematic of a committed bracket (illustrative)
\begin{scope}[shift={(9.55,1.15)}, scale=1.05]
  \filldraw[fill=ezaccentsoft, draw=ezink, line width=0.6pt]
    (0,0) -- (3.6,0) -- (3.6,0.32) -- (1.25,0.32) .. controls (0.95,0.95) .. (0.62,1.55) -- (0.62,2.55) -- (0,2.55) -- cycle;
  \filldraw[fill=ezaccent!25, draw=ezink, line width=0.5pt] (0.62,0.32) -- (1.15,0.32) .. controls (0.88,0.85) .. (0.62,1.25) -- cycle;
  \draw[ezink, line width=0.5pt] (0.62,0.32) -- (0.62,1.55);
  \foreach \x in {2.3,3.15} \filldraw[fill=white, draw=ezink, line width=0.4pt] (\x,0.16) circle (0.08);
  \filldraw[fill=white, draw=ezink, line width=0.4pt] (0.31,2.2) circle (0.09);
  \fill[pattern=north east lines, pattern color=ezgrey] (-0.1,-0.16) rectangle (3.7,-0.03);
  \draw[ezwarm, line width=1pt, -{Stealth[length=2.2mm]}] (1.35,2.95) -- (0.4,2.45);
  \node[font=\sffamily\scriptsize, text=ezwarm] at (1.55,3.05) {$F$};
  \node[font=\sffamily\fontsize{5.8pt}{6.4pt}\selectfont, text=ezgrey, align=center, anchor=north] at (1.8,-0.25)
    {committed part (schematic)\\offset web and gusset};
\end{scope}
\synfigtag
\end{tikzpicture}

\caption{CAD parts. Left: mass against compliance, both relative, for one committed proposal per
specification; lower and to the left is better, and points right of the dashed line violate the
stiffness limit. Right: schematic of a committed bracket with an offset web and gusset, the structural
motif examined in dossier D-C1 (\cref{sec:discoveries}). Synthetic (gaps EZ-G17, EZ-G19).}
\label{fig:cad}
\end{figure}

\subsection{Against frontier models at one attempt}

Frontier language models write strong programs from descriptions: in the program's own pilot, a
general regressor written by one without seeing any data ranked first of eighteen methods \pre{}
(\cref{app:history}). Given the byte-identical information packet, no tools and one attempt, the best of
three frontier models trails EurekaZero by \syn[EZ-G07]{\EZfrA}, \syn[EZ-G07]{\EZfrB} and
\syn[EZ-G07]{\EZfrC} utility in the three games (\cref{tab:frontier}). The gap is largest for CAD,
where only \syn[EZ-G07]{\EZfrValidC\%} of frontier construction programs built valid solids at the
first attempt; text-to-CAD systems usually repair failures with retries \citep{wang2026text2cadbench},
which the Zero contract forbids and which we report separately as search.

\begin{table}[t]
\caption{One attempt with matched information: EurekaZero against three frontier models. Synthetic
(gap EZ-G07).}
\label{tab:frontier}
\centering\tablesf
\begin{tabular}{@{}lrrrrrr@{}}
\toprule
 & \multicolumn{4}{c}{Utility of the first proposal} & \multicolumn{2}{c}{Valid first proposals (\%)} \\
\cmidrule(lr){2-5}\cmidrule(l){6-7}
Game & EurekaZero & Frontier 1 & Frontier 2 & Frontier 3 & EurekaZero & Best frontier \\
\midrule
A Learning procedures & \syn[EZ-G07]{\EZuA} & \syn[EZ-G07]{\EZfrOneA} & \syn[EZ-G07]{\EZfrTwoA} & \syn[EZ-G07]{\EZfrThreeA} & \syn[EZ-G07]{\EZvalidA} & \syn[EZ-G07]{\EZfrValidA} \\
B1 Next experiment & \syn[EZ-G07]{\EZuB} & \syn[EZ-G07]{\EZfrOneB} & \syn[EZ-G07]{\EZfrTwoB} & \syn[EZ-G07]{\EZfrThreeB} & \syn[EZ-G07]{\EZvalidB} & \syn[EZ-G07]{\EZfrValidB} \\
C CAD parts & \syn[EZ-G07]{\EZuC} & \syn[EZ-G07]{\EZfrOneC} & \syn[EZ-G07]{\EZfrTwoC} & \syn[EZ-G07]{\EZfrThreeC} & \syn[EZ-G07]{\EZbuildC} & \syn[EZ-G07]{\EZfrValidC} \\
\bottomrule
\end{tabular}
\eztablenote{Frontier models are ranked by mean utility within each game;
model identities, versions and access dates are recorded with the confirmation roster.}
\end{table}

\subsection{Where the gains disappear}
\label{sec:heterogeneity}

The predeclared subgroup analyses locate the limits of the first proposal. In learning procedures the
benefit shrinks with the number of training rows and is indistinguishable from zero above 50,000
rows \synbadge. In next-experiment decisions it disappears above 20 dimensions, where every method's
first query is close to uninformed \synbadge. In CAD it falls for specifications with more than three
load cases or with contact, which the bounded language represents poorly, and failures concentrate
there \synbadge. These are the regimes where search keeps its advantage and where a first proposal is
best used as a seed (\cref{sec:compute}).
